% =========================================================
% Community-Aware Early Warning of Distrust Emergence
% in Temporal Bitcoin Trust Networks
%
% IEEE conference paper draft.
% Compile with pdflatex; uses IEEEtran.cls (IEEEtran v1.8b or later).
%
% This is the integrated final version: rolling-origin and paired
% statistical testing are now the primary evidence; the locked
% single-fold analysis (bootstrap, ROC/PR curve overlays, feature
% importance) is preserved as Section~\ref{sec:results:locked}
% (secondary). A compact comparison with heuristic baselines and
% LightGBM is added as Section~\ref{sec:results:baselines}.
% =========================================================
\documentclass[conference]{IEEEtran}

\IEEEoverridecommandlockouts
% packages
\usepackage{cite}
\usepackage{amsmath,amssymb,amsfonts}
\usepackage{graphicx}
\usepackage{textcomp}
\usepackage{booktabs}
\usepackage{url}
\usepackage{hyperref}
\hypersetup{hidelinks}

% Avoid IEEEtran title-style font conflicts
\def\BibTeX{{\rm B\kern-.05em{\sc i\kern-.025em b}\kern-.08em
    T\kern-.1667em\lower.7ex\hbox{E}\kern-.125emX}}

% Inline code shortcut
\newcommand{\code}[1]{\texttt{#1}}


\begin{document}

% =========================================================
% TITLE BLOCK
% =========================================================
\title{Community-Aware Early Warning of Distrust Emergence \\
in Temporal Bitcoin Trust Networks}

\author{
\IEEEauthorblockN{\"Omer Faruk G\"uler}
\IEEEauthorblockA{
Hacettepe University\\
Ankara, Turkey\\
omer55glr@gmail.com
}
}

\maketitle

% =========================================================
% ABSTRACT
% =========================================================
\begin{abstract}
Bitcoin-OTC was a peer-rated trust marketplace whose users rated each
other on a $[-10, +10]$ scale, producing a signed, weighted, directed,
timestamped social network. We study early warning of \emph{new
negative} trust edges: at the end of month $t$, can we rank ordered
pairs of users by their risk of receiving a new negative directed
rating in month $t{+}1$, using only information available at time $t$?
We build an interpretable, leakage-controlled pipeline with cumulative
monthly snapshots and a three-way feature ablation
(structural; structural + temporal; structural + temporal + community-aware)
under Logistic Regression and Random Forest, with Louvain communities
extracted per snapshot. We evaluate under a rolling-origin protocol
over 25 disjoint single-month test folds covering 425 total positive
events. Adding temporal features yields a statistically significant
ROC-AUC improvement under both models (paired Wilcoxon
$p \leq 1.4\!\times\!10^{-4}$ after Holm-Bonferroni correction on
four pre-declared primary tests; matched-pairs rank-biserial
$r_{\mathrm{rb}} \approx +0.8$), and a significant Random Forest PR-AUC
improvement ($p = 6.3\!\times\!10^{-4}$, $r_{\mathrm{rb}} = +0.74$). A
parameter-free recency-only ranker already reaches ROC-AUC $0.886$ on
the same protocol -- comparable to structural-only learned models --
so a substantial fraction of the temporal gain is attributable to
recency alone; the engineered temporal family adds a real but modest
increment above pure recency. Adding community-aware features on top
of temporal features yields a small, directionally consistent Random
Forest PR-AUC gain (19 of 25 folds positive; sign test $p = 0.015$;
Wilcoxon $p = 0.048$, marginal under the pre-declared threshold) and
does not improve ROC-AUC ($p = 0.096$); under Logistic Regression the
community effect is undetectable on any metric, and under LightGBM
the community PR-AUC effect is not detectable
(mean $-0.007$, $p = 0.51$). We therefore frame temporal lift as
the robust headline finding and community lift as a small,
model-dependent refinement under Random Forest. Precision@$k$ proves
unstable under the rolling protocol -- per-fold median P@50 is zero --
and is reported as a supplementary metric.
\end{abstract}

\begin{IEEEkeywords}
signed networks, temporal link prediction, community detection,
Louvain, Random Forest, rolling-origin evaluation, Wilcoxon
signed-rank test, Bitcoin-OTC, weak ties.
\end{IEEEkeywords}

% =========================================================
\section{Introduction}
% =========================================================

Online trust marketplaces let participants rate one another so that
newcomers can assess the reliability of potential counterparties.
Bitcoin-OTC, an over-the-counter cryptocurrency marketplace, hosted
such a web-of-trust between 2010 and early 2016. Each rating is a
directed signed weighted timestamped event: user $u$ assigns user $v$
an integer rating in $[-10, +10]$, and the rating is stored with its
timestamp. The resulting graph is one of the canonical benchmarks for
signed network analysis~\cite{kumar2016}.

This paper studies \emph{early warning of distrust emergence}: given
the network at the end of month $t$, we want to rank ordered pairs
$(u, v)$ by the probability that $v$ will receive a new \emph{negative}
directed rating from $u$ in month $t+1$. Useful applications include
marketplace moderation, fraud surveillance, and reputation governance.
The task differs from generic link prediction in three ways. First,
negative ratings are rare: only about 10\% of edges in Bitcoin-OTC are
negative, and new negative events per month are far rarer still.
Second, the task is sign-dependent: predicting any new edge and
predicting a new negative edge are different problems, and positive-edge
popularity heuristics do not directly transfer. Third, the evaluation
must respect the arrow of time -- a model that uses post-$t$
information would not be deployable for warning.

Our research question is:
\begin{quote}
At snapshot $t$, can we predict whether an ordered pair $(u, v)$ --
either previously unconnected, or connected only positively -- will
receive a new negative directed edge in the window
$(t, t{+}1]$, using only information available at time $t$?
\end{quote}

Our contributions are:
\begin{itemize}
\item An interpretable feature-based pipeline -- classical
link-prediction, signed structural, trust-summary, temporal, and
community-aware features -- with explicit leakage controls and a
time-based candidate generation procedure.
\item A three-way feature ablation (structural; structural + temporal;
structural + temporal + community) under Logistic Regression and
Random Forest.
\item A \emph{rolling-origin evaluation} over 25 disjoint single-month
test folds (425 total positive events), with paired per-fold metric
deltas across feature sets.
\item Formal paired statistical testing on per-fold deltas: two-sided
Wilcoxon signed-rank and exact binomial sign tests, with
Holm-Bonferroni correction on a pre-declared primary family of four
tests and Benjamini-Hochberg FDR control on twenty secondary tests.
\item A comparison of the feature-based pipeline against four
parameter-free heuristic baselines and LightGBM, evaluated under the
same rolling-origin protocol, to assess how much of the temporal
gain is captured by simple recency information and whether a stronger
tabular learner improves the picture.
\item A descriptive analysis of community structure and weak-tie
signals, separating global feature-level evidence from a small local
top-of-ranked-list diagnostic.
\end{itemize}

The pipeline is intentionally lighter than the deep signed-graph
methods surveyed by Zhang \emph{et al.}~\cite{zhang2024}. Our goal is
interpretable analysis on a single dataset; we do not claim
state-of-the-art absolute performance.

% =========================================================
\section{Related Work}
% =========================================================

\textbf{Signed networks and weighted trust prediction.} The
Bitcoin-OTC and Bitcoin-Alpha datasets, together with the iterative
\emph{fairness} and \emph{goodness} statistics for nodes, were
introduced by Kumar \emph{et al.}~\cite{kumar2016}. We do not
reimplement iterated fairness/goodness here; instead, we use four
lightweight per-node trust-summary features (mean rating given, mean
rating received, negative-given ratio, negative-received ratio) as a
deliberately simple analogue.

\textbf{Behavioral dynamics on Bitcoin-OTC.} Bertazzi
\emph{et al.}~\cite{bertazzi2018} analyze the multilayer structure of
the Bitcoin-OTC web of trust, contrasting rewarding and punitive
behaviors at the network level. This motivates our temporal framing
and the inclusion of recency-based features (days since last activity,
recent-window counts, and exponentially decayed aggregates).

\textbf{Community-aware temporal prediction.}
Choudhury~\cite{choudhury2024} develops community-aware dynamic
similarity features for supervised link prediction on dynamic
networks, arguing that community structure carries predictive signal
beyond pair-level topology. We operationalize this idea by running
Louvain on the positive-edge projection of every snapshot and
engineering community-aware pair features
(\code{same\_community}, community sizes, \code{neighborhood\_overlap},
and boundary fractions).

\textbf{Deep signed graph learning.} Zhang \emph{et al.}~\cite{zhang2024}
survey signed graph representation learning, including
GNN-based approaches such as signed graph convolutional and attention
networks. We use this body of work for positioning only: our approach
is shallow and interpretable, and we do not benchmark against GNNs.

\textbf{Classical heuristics and structural theory.} We use Common
Neighbors, Jaccard similarity, Adamic-Adar~\cite{adamic2003}, and
Preferential Attachment on the undirected any-sign projection of the
snapshot graph. Our community-aware features include the
neighborhood-overlap diagnostic associated with Granovetter's
strength-of-weak-ties hypothesis~\cite{granovetter1973}. Louvain
community detection follows Blondel \emph{et al.}~\cite{blondel2008}.

% =========================================================
\section{Dataset and Problem Formulation}
% =========================================================

\subsection{Bitcoin-OTC and basic terminology}
Each row of the public dataset (file \code{soc-sign-bitcoinotc.csv},
SNAP~\cite{snap}) is a single rating event with fields
\code{(source, target, rating, timestamp)}. A \emph{node} is a user.
A \emph{directed edge} $u \to v$ is a rating event from $u$ to $v$. The
\emph{sign} of an edge is $+1$ if $\text{rating} > 0$ and $-1$ if
$\text{rating} < 0$. The \emph{weight} of an edge is the integer
rating in $[-10, +10]$. The \emph{timestamp} is a Unix-seconds float;
the dataset spans 2010-11-08 to 2016-01-25, approximately 63 calendar
months.

After defensive cleaning (drop self-loops and zero ratings; neither
occurs in this file), the corpus has 35{,}592 rating events, 5{,}881
distinct users, 32{,}029 positive and 3{,}563 negative ratings
(approximately 9:1 imbalance), and reciprocity $0.79$. The giant
weakly connected component has 5{,}875 nodes and the giant strongly
connected component has 4{,}709 nodes. Fig.~\ref{fig:edges_per_month}
shows the monthly interaction volume; Fig.~\ref{fig:rating_histogram}
shows the rating distribution.

\begin{figure}[t]
\centering
\includegraphics[width=\columnwidth]{results/figures/edges_per_month.png}
\caption{Monthly interaction volume on Bitcoin-OTC, Nov-2010 through
Jan-2016. Activity grows through 2014--2015 and then collapses; the
tail collapse motivates the eligibility filter discussed in
Section~\ref{sec:setup}.}
\label{fig:edges_per_month}
\end{figure}

\begin{figure}[t]
\centering
\includegraphics[width=\columnwidth]{results/figures/rating_histogram.png}
\caption{Distribution of rating values across all events. Positive
ratings (especially $+1$ and $+10$) dominate; negatives are a long left
tail.}
\label{fig:rating_histogram}
\end{figure}

\subsection{Cumulative monthly snapshots}
We adopt a cumulative monthly schedule. For each month index $t \in
\{0, \dots, 62\}$, the snapshot graph at time $t$ contains all events
with timestamp at or before the end of month $t$. When the pair
$(u, v)$ appears multiple times in the cumulative window, we collapse
the multi-edge to its \emph{latest} rating before the end of month
$t$. This encodes ``$u$'s current trust judgment of $v$'' and is used
\emph{only for the snapshot graph state}. Temporal and history-based
features instead use the full event history with
$\mathrm{ts} \leq \mathrm{end\_of\_month}(t)$, preserving the
underlying multi-event structure.

From the directed snapshot graph $G^t_{\mathrm{dir}}$ we derive two
undirected weighted projections: $G^t_{\mathrm{pos}}$ keeps only edges
whose collapsed sign is $+1$, with weight equal to the number of
distinct directed positive contributions; $G^t_{\mathrm{neg}}$ is the
analogous negative projection.

\subsection{The task}
Given snapshot $t$, the target is a binary label per ordered candidate
pair $(u, v)$:
\begin{equation*}
y(u, v) =
\begin{cases}
1, & \exists\,(u, v, r, \mathrm{ts}) \text{ with } r < 0 \\
   & \text{and } \mathrm{end\_of\_month}(t) < \mathrm{ts} \\
   & \quad\leq \mathrm{end\_of\_month}(t{+}1), \\
0, & \text{otherwise.}
\end{cases}
\end{equation*}
Both ``no new $u \to v$ event'' and ``only new positive $u \to v$
events'' map to $y = 0$; only a new \emph{negative} directed rating
counts. This includes the case where $u$ previously rated $v$
positively but flips to negative in the next window, which is a
legitimate form of distrust emergence and is therefore retained as a
positive example.

\subsection{Why ranking metrics rather than thresholded
precision/recall}
The candidate set for any non-trivial snapshot is large (approximately
$1 \times 10^5$ to $2 \times 10^5$ ordered pairs after filtering;
see Section~\ref{sec:method}), and the prior of $y = 1$ is on the
order of $10^{-4}$. At this imbalance, a threshold-$0.5$ classifier is
not a meaningful operating point. A model with
\code{class\_weight = 'balanced'} tends to predict the positive class
at unrealistic precision; a default Random Forest predicts essentially
no positives. Threshold-dependent precision, recall, and F1 are
therefore nearly uninformative in this regime. We focus on
\emph{ranking} metrics: ROC-AUC, PR-AUC (average precision), and
Precision@$k$ for small $k$. Precision@$k$ proves unstable across
folds and is reported as supplementary
(Section~\ref{sec:results:community}); ROC-AUC and PR-AUC carry the
inferential weight.

% =========================================================
\section{Methodology}\label{sec:method}
% =========================================================

\subsection{Core node filtering}
A node $n$ is \emph{core at snapshot $t$} iff (a) its total
interactions (in + out events) up to $\mathrm{end\_of\_month}(t)$ is at
least \code{MIN\_ACTIVITY = 3}, and (b) its most recent activity is
within \code{RECENT\_DAYS = 180} of $\mathrm{end\_of\_month}(t)$. Both
thresholds are configuration parameters; we adopt these values as
starting defaults and do not sweep them.

\subsection{Candidate pair generation}
For each snapshot $t$, candidate ordered pairs $(u, v)$ satisfy:
(1)~$u, v \in \mathrm{Core}(t)$ and $u \neq v$;
(2)~$v$ is within \code{K\_HOP = 2} hops of $u$ on the undirected
any-sign projection of $G^t_{\mathrm{dir}}$;
(3)~there is no existing \emph{negative} directed edge $u \to v$ at
time $t$ -- pairs with a prior \emph{positive} $u \to v$ are retained,
because a rating flip from positive to negative in the next window is
itself distrust emergence;
(4)~the total candidate count is capped at
\code{MAX\_CANDIDATES = 200{,}000} via fixed-seed uniform
downsampling. The cap binds only at the densest mid-period snapshots.

\subsection{Label generation and leakage control}
Labels are derived from the raw event history in the window
$(\mathrm{end\_of\_month}(t), \mathrm{end\_of\_month}(t{+}1)]$. A
defensive assertion in every feature precompute confirms that the
maximum timestamp used during feature computation does not exceed
$\mathrm{end\_of\_month}(t)$.

\subsection{Feature families}
The full feature vector has 36 columns split across five families:

\textbf{Classical link prediction (4 features).} Common Neighbors,
Jaccard similarity, Adamic-Adar~\cite{adamic2003}, Preferential
Attachment, all computed on the undirected any-sign projection of
$G^t_{\mathrm{dir}}$.

\textbf{Signed structural (8 features).} Positive and negative
out-degree of $u$; positive and negative in-degree of $v$; positive
and negative common-neighbor counts (on $G^t_{\mathrm{pos}}$ and
$G^t_{\mathrm{neg}}$ respectively); reciprocity flag (1 if $v \to u$
exists in $G^t_{\mathrm{dir}}$); prior-interaction flag (1 if
$u \to v$ exists in $G^t_{\mathrm{dir}}$ -- necessarily a positive
prior, by construction of the candidate set, so this flag is partly
degenerate and we report it transparently).

\textbf{Trust summary (4 features).} Mean rating given by $u$; mean
rating received by $v$; fraction of $u$'s given ratings that were
negative; fraction of $v$'s received ratings that were negative.
These are a deliberately simple, non-iterated analogue of the
fairness/goodness statistics in~\cite{kumar2016}.

\textbf{Temporal / recency (12 features).} Days since $u$'s last
activity; days since $v$'s last activity; days since the last
$u \to v$ interaction (sentinel value $9999$ when no prior $u \to v$
event exists); count of past $u \to v$ events; recent-window (90-day)
counts of negatives given by $u$, negatives received by $v$, total
events given by $u$, total events received by $v$; and four
exponentially decayed aggregates with half-life $\tau = 60$ days
(decayed total activity and decayed negative count, for both $u$ and
$v$).

\textbf{Community-aware (8 features).} \code{same\_community}
indicator; its complement \code{cross\_community\_pair};
\code{community\_size\_u}, \code{community\_size\_v}; the log-size
gap $|\log s_u - \log s_v|$;
\code{neighborhood\_overlap} $= |N(u) \cap N(v) \setminus \{u,v\}|
/ |N(u) \cup N(v) \setminus \{u,v\}|$ on the undirected any-sign
projection; and \code{boundary\_frac\_u}, \code{boundary\_frac\_v},
each defined for a node $n$ as the fraction of $n$'s neighbors in
$G^t_{\mathrm{pos}}$ whose Louvain community differs from $n$'s.

\subsection{Community extraction}\label{sec:method:community}
For each snapshot $t$ we run Louvain~\cite{blondel2008} on the
positive-edge undirected weighted graph $G^t_{\mathrm{pos}}$. Louvain
is stochastic; we fix \code{random\_state} and cache the resulting
partition per snapshot so that downstream community-feature
extraction is deterministic. Across the 37 split snapshots (see
Section~\ref{sec:setup}) modularity sits in $[0.47, 0.52]$ with a slight
upward trend, while the number of communities grows from $\sim 17$ at
$t = 13$ to $\sim 40$--$47$ at $t = 49$ (Fig.~\ref{fig:modularity}).

\begin{figure}[t]
\centering
\includegraphics[width=\columnwidth]{results/figures/modularity_over_time.png}
\caption{Louvain modularity (left axis) and number of communities
(right axis) on the per-snapshot positive-edge graph
$G^t_{\mathrm{pos}}$. Modularity is stable in $[0.47, 0.52]$; the
community count grows roughly linearly as the network grows. Dotted
lines mark the train\,/\,val and val\,/\,test boundaries of the
locked single-fold analysis (Section~\ref{sec:results:locked}).}
\label{fig:modularity}
\end{figure}

\subsection{Models}
We train three models per feature set: (i) a \emph{random baseline}
that draws uniform probabilities (seeded), to calibrate the floor on
each metric; (ii) \emph{Logistic Regression} with \code{StandardScaler}
and \code{class\_weight = 'balanced'}, \code{liblinear} solver,
\code{max\_iter = 2000}; (iii) \emph{Random Forest} with
\code{n\_estimators = 200}, \code{min\_samples\_leaf = 20},
\code{class\_weight = 'balanced\_subsample'}. We perform no
hyperparameter tuning beyond these defaults; the focus is feature
engineering and ablation. Section~\ref{sec:results:baselines}
additionally evaluates four parameter-free heuristic baselines and a
fixed-configuration LightGBM as a stronger learner.

\subsection{Training-set negative subsampling}
The training data for any rolling-origin fold (Section~\ref{sec:setup})
contains on the order of $10^{6}$ candidate pairs and only a few
hundred positives. To keep training tractable while preserving all
positives, we subsample negatives in TRAIN to a $100{:}1$
negative:positive ratio. Validation and test sets remain complete;
this is a training acceleration, not an evaluation trick. A fixed
RNG seed makes the subsample identical across feature sets and models
within a fold, so per-fold paired comparisons between feature sets
within a model are valid.

\subsection{Evaluation metrics}
We report ROC-AUC, PR-AUC (average precision), and Precision@$k$ for
$k \in \{50, 100\}$. Threshold metrics (precision, recall, F1 at the
default 0.5 cutoff) are recorded but read as uninformative at the
present imbalance; we do not use them for inference.

% =========================================================
\section{Experimental Setup}\label{sec:setup}
% =========================================================

\subsection{Eligible snapshot pool}
A snapshot is \emph{eligible} for evaluation if (i) its next-month
window contains at least \code{MIN\_POS\_FOR\_SPLIT = 10} raw new
negative events, and (ii) it sits in the longest contiguous block of
eligible snapshots (this excludes a sporadic late-dataset island at
$t = 53$). The resulting eligible pool is $t \in \{13, \dots, 49\}$
(37 snapshots).

\subsection{Rolling-origin evaluation protocol (primary)}
\label{sec:setup:rolling}

Our primary evaluation is an expanding-window rolling-origin protocol
over the eligible pool. For each fold index
$k \in \{25, 26, \dots, 49\}$ (25 disjoint single-month test folds):
\begin{itemize}
\item TRAIN  = candidate pairs from snapshots $\{13, \dots, k{-}2\}$,
\item VAL    = candidate pairs from snapshot $\{k{-}1\}$ (held out; not used for tuning),
\item TEST   = candidate pairs from snapshot $\{k\}$.
\end{itemize}
We use an expanding-window TRAIN rather than a sliding window because
positive events are rare and discarding early ones would weaken
training. We use single-month TEST windows rather than two-month
windows because they maximize the number of disjoint folds (25 vs
$\sim$12), match an operational ``next-month watchlist'' deployment
cadence, and produce fold-level metric vectors directly usable for
paired non-parametric tests. The total number of positive test events
across all folds is $\sum_{k} n_{\mathrm{pos}}(k) = 425$, an
approximately $12\!\times$ increase over the single locked fold
described in Section~\ref{sec:setup:locked}.

For each fold and each (model, feature set) combination, we fit a
fresh model on the fold's TRAIN with the same $100{:}1$ training
negative subsampling as in Section~\ref{sec:method}, score the fold's
TEST, and record ROC-AUC, PR-AUC, Precision@50, and Precision@100.
All seeds are fixed; within a fold the training subsample is
identical across (model, feature set), so per-fold deltas across
feature sets are properly paired.

\subsection{Paired statistical testing on per-fold deltas}
\label{sec:setup:stats}

On the $n = 25$ per-fold deltas we compute, per (model, comparison,
metric):
\begin{itemize}
\item a two-sided paired Wilcoxon signed-rank test
(\code{zero\_method='wilcox'});
\item a two-sided exact binomial sign test on the direction counts;
\item the matched-pairs rank-biserial correlation
$r_{\mathrm{rb}} = (W_+ - W_-)/(W_+ + W_-)$ as a Wilcoxon effect
size.
\end{itemize}

We declare a small \emph{primary family} of four tests in advance --
the ones we consider central to the paper's question: RF (S+T) - S on
ROC-AUC, RF (S+T) - S on PR-AUC, RF (S+T+C) - (S+T) on PR-AUC, and
LR (S+T) - S on ROC-AUC. We control family-wise error at
$\alpha = 0.05$ on this primary set with Holm-Bonferroni. All other
(model, comparison, metric) combinations form a secondary family, on
which we control the false-discovery rate at $q = 0.05$ with
Benjamini-Hochberg. We report raw p-values, the corrected reject
decision, and the rank-biserial effect size for every test. When the
per-fold delta is exactly zero on many folds (particularly for
Precision@$k$, where the discrete probability grid produces many
ties), the Wilcoxon test becomes unreliable and we rely on the sign
test only.

\subsection{Held-out single-fold analysis (secondary)}
\label{sec:setup:locked}

For continuity with prior pipeline stages and to support the bootstrap
analysis in Section~\ref{sec:results:bootstrap}, we also retain the
last two eligible snapshots ($t \in \{48, 49\}$) as a single locked
held-out fold, with the immediately preceding two snapshots
($t \in \{46, 47\}$) as a tuning-free validation set
(Table~\ref{tab:split}). All numbers from this single fold are
reported in Section~\ref{sec:results:locked}; they are kept as a
secondary analysis. The headline numerical claims in this paper come
from the rolling-origin protocol.

\begin{table}[t]
\caption{Locked single held-out fold. Used in
Section~\ref{sec:results:locked} for the bootstrap analysis, ROC/PR
curve overlays, and feature importance; the headline numerical
claims come from the rolling-origin protocol of
Section~\ref{sec:setup:rolling}.}
\label{tab:split}
\centering
\begin{tabular}{lccc}
\toprule
Split & Snapshots & Pairs & Positives \\
\midrule
TRAIN & $t = 13 \dots 45$ & 5{,}013{,}346 & 545 \\
VAL   & $t = 46, 47$      &   184{,}648   &  16 \\
TEST  & $t = 48, 49$      &   148{,}639   &  34 \\
\bottomrule
\end{tabular}
\end{table}

\subsection{Bootstrap protocol (for the locked single fold)}
\label{sec:setup:bootstrap}
For the locked single-fold analysis we attach uncertainty estimates
with a paired stratified bootstrap. For $B = 1000$ iterations, we
draw 34 positive indices with replacement and 148{,}605 negative
indices with replacement (preserving the original prevalence and
guaranteeing both labels in every resample). The same resample index
vector is reused across all three feature sets within one model, so
deltas such as
$(\mathrm{S}{+}\mathrm{T}{+}\mathrm{C}) - (\mathrm{S}{+}\mathrm{T})$
are properly paired. We add a $\mathcal{U}[0, 10^{-12}]$ jitter to RF
probabilities before top-$k$ ranking to randomize ties on the
discrete $1/200$ leaf-vote grid without perturbing non-tied rankings.
We report 95\% percentile confidence intervals and the
bootstrap-positive share $\mathrm{frac\_boot}(>0)$; this share is
\emph{not} a frequentist p-value, and we interpret a CI that overlaps
zero as ``not distinguishable from sampling noise on this test set.''
This analysis quantifies uncertainty conditional on the locked fold;
the rolling-origin protocol of Section~\ref{sec:setup:rolling}
provides the cross-fold significance picture.

\subsection{Reproducibility and artifacts}
All seeds (model fits, training subsample, Louvain partition,
bootstrap) are fixed at $42$. Intermediate parquet files and the
locked-fold test-probability matrix
(\code{results/test\_predictions.parquet}) are persisted so that
downstream analysis re-runs without re-training. Per-fold rolling-origin
probabilities for the Random Forest S$+$T$+$C model are stored under
\code{results/rolling\_predictions/}. Code is organized in
\code{src/} and \code{scripts/}; the headline numbers in
Section~\ref{sec:results} come from
\code{results/rolling\_eval\_per\_fold.csv} (rolling-origin),
\code{results/rolling\_stats\_tests.csv} (paired tests),
\code{results/stronger\_baselines\_per\_fold.csv} (heuristics and
LightGBM), and
\code{results/test\_metrics\_from\_predictions.csv} (locked fold).

% =========================================================
\section{Results}\label{sec:results}
% =========================================================

We label feature sets S (structural + trust summary, 16 columns),
S$+$T (+temporal, 28 columns), and S$+$T$+$C (+community, 36 columns).

\subsection{Rolling-origin metrics}\label{sec:results:rolling}

Table~\ref{tab:rolling_summary} reports across-fold mean and standard
deviation per (model, feature set) over the $n = 25$ rolling-origin
folds (total 425 positive test events). Adding temporal features
clearly improves both ROC-AUC and PR-AUC under both models; adding
community-aware features on top of temporal moves PR-AUC slightly
upward under Random Forest while leaving Logistic Regression
essentially unchanged.

\begin{table}[t]
\caption{Rolling-origin metrics across 25 single-month test folds
(total 425 positives). Values are across-fold mean $\pm$ std.}
\label{tab:rolling_summary}
\centering
\small
\begin{tabular}{llcccc}
\toprule
F.\,set & Model & ROC-AUC & PR-AUC & P@50 & P@100 \\
\midrule
S       & lr & $0.889 \pm 0.077$ & $0.006 \pm 0.006$ & $0.012$ & $0.008$ \\
S$+$T   & lr & $0.930 \pm 0.046$ & $0.007 \pm 0.008$ & $0.014$ & $0.010$ \\
S$+$T$+$C & lr & $0.928 \pm 0.048$ & $0.007 \pm 0.009$ & $0.014$ & $0.010$ \\
\addlinespace
S       & rf & $0.881 \pm 0.073$ & $0.006 \pm 0.011$ & $0.009$ & $0.008$ \\
S$+$T   & rf & $0.924 \pm 0.054$ & $0.014 \pm 0.023$ & $0.022$ & $0.016$ \\
S$+$T$+$C & rf & $0.920 \pm 0.058$ & $0.017 \pm 0.029$ & $0.022$ & $0.016$ \\
\bottomrule
\end{tabular}
\end{table}

\subsection{Paired statistical tests}\label{sec:results:tests}

Table~\ref{tab:primary_tests} reports the four pre-declared primary
tests on the per-fold deltas, with Holm-Bonferroni decisions at
$\alpha = 0.05$. All four nulls are rejected after correction, but
the RF community PR-AUC test is marginal -- its Wilcoxon $p = 0.048$
sits at the loosest Holm threshold. The sign-test p-value
($p = 0.015$) on the same comparison is more robust to ties.
Fig.~\ref{fig:paired_deltas} visualizes the per-fold deltas for
these four comparisons.

\begin{table}[t]
\caption{Primary paired statistical tests on rolling-origin per-fold
deltas ($n = 25$; two-sided paired Wilcoxon signed-rank;
Holm-Bonferroni at $\alpha = 0.05$). $r_{\mathrm{rb}}$ is the
matched-pairs rank-biserial effect size.}
\label{tab:primary_tests}
\centering
\small
\begin{tabular}{llccc}
\toprule
Comparison & Metric & Wilcoxon $p$ & sign $p$ & $r_{\mathrm{rb}}$ \\
\midrule
RF $(\mathrm{S}{+}\mathrm{T}) - \mathrm{S}$       & ROC-AUC & $4.5{\times}10^{-5}^{\dagger}$ & $1.6{\times}10^{-4}$ & $+0.85$ \\
LR $(\mathrm{S}{+}\mathrm{T}) - \mathrm{S}$       & ROC-AUC & $1.4{\times}10^{-4}^{\dagger}$ & $9.1{\times}10^{-4}$ & $+0.81$ \\
RF $(\mathrm{S}{+}\mathrm{T}) - \mathrm{S}$       & PR-AUC  & $6.3{\times}10^{-4}^{\dagger}$ & $9.1{\times}10^{-4}$ & $+0.74$ \\
RF $(\mathrm{S}{+}\mathrm{T}{+}\mathrm{C}) - (\mathrm{S}{+}\mathrm{T})$ & PR-AUC  & $4.8{\times}10^{-2}^{\dagger}$ & $1.5{\times}10^{-2}$ & $+0.45$ \\
\bottomrule
\end{tabular}\\[2pt]
\raggedright\footnotesize${}^{\dagger}$ rejects $H_0$ after Holm-Bonferroni at $\alpha = 0.05$.
\end{table}

\begin{figure}[t]
\centering
\includegraphics[width=\columnwidth]{results/figures/rolling_paired_deltas.png}
\caption{Per-fold paired deltas for the four primary comparisons
(blue: positive fold; red: negative fold).}
\label{fig:paired_deltas}
\end{figure}

Among secondary tests (Benjamini-Hochberg FDR $q = 0.05$): the full
$\mathrm{S} \to \mathrm{S}{+}\mathrm{T}{+}\mathrm{C}$ delta is
significant on both metrics and both models;
$\mathrm{LR}\,(\mathrm{S}{+}\mathrm{T}) - \mathrm{S}$ on PR-AUC is
significant despite a small absolute magnitude;
$\mathrm{RF}\,(\mathrm{S}{+}\mathrm{T}) - \mathrm{S}$ on
Precision@100 is significant. No community-vs-temporal ROC-AUC test
reaches significance: RF (S+T+C) - (S+T) ROC-AUC has $p = 0.096$ with
only 10/25 folds positive, consistent with saturation; LR (S+T+C) -
(S+T) ROC-AUC has $p = 0.47$. LR community-vs-temporal PR-AUC is
also null ($p = 0.35$). Most Precision@$k$ tests either fail to
reach significance after correction or are flagged unreliable
because of tied zeros.

\subsection{Temporal ablation}\label{sec:results:temporal}

The temporal-over-structural ROC-AUC delta is positive in 22 of 25
RF folds (median $+0.031$, mean $+0.044$) and in 21 of 25 LR folds
(median $+0.036$, mean $+0.041$). Both pass Holm correction with
very small p-values and large effect sizes
($r_{\mathrm{rb}} \approx +0.85$ and $+0.81$).

The Random Forest temporal PR-AUC delta is also significant after
Holm correction (Wilcoxon $p = 6.3\!\times\!10^{-4}$, sign-test
$p = 9.1\!\times\!10^{-4}$, $r_{\mathrm{rb}} = +0.74$), with 21 of
25 folds positive. The Logistic Regression temporal PR-AUC delta is
significant under the secondary BH correction
($p = 3.4\!\times\!10^{-3}$), but its mean magnitude ($+0.002$) is
small in absolute terms.

The temporal feature family is internally redundant by design:
recent-window counts and exponentially decayed aggregates carry
overlapping information, with within-pair correlations $\geq 0.9$
for some pairs. We retain both for transparency. The qualitative
weight of this lift is examined further in
Section~\ref{sec:results:baselines}, where a parameter-free
recency-only ranker turns out to capture a substantial fraction of
it.

\subsection{Community-aware ablation}\label{sec:results:community}

\textbf{Community PR-AUC under Random Forest is small but
directionally consistent.} The (S+T+C) - (S+T) PR-AUC delta is
positive in 19 of 25 folds, with median $+0.001$ and mean $+0.003$.
The Wilcoxon test sits right at the Holm threshold
($p = 0.048$); the sign test is more robust ($p = 0.015$). The
rank-biserial effect size $r_{\mathrm{rb}} = +0.45$ is in the medium
range. We describe this gain as small but directionally consistent
across folds, with marginal statistical support. The effect is also
model-specific: under LightGBM the same comparison has a slightly
negative point estimate (Section~\ref{sec:results:baselines}).

\textbf{Community does not improve ROC-AUC over temporal.} The
(S+T+C) - (S+T) ROC-AUC delta is positive in only 10 of 25 RF folds
($p = 0.096$, $r_{\mathrm{rb}} = -0.39$) and 12 of 25 LR folds
($p = 0.47$). The Random Forest direction is mildly negative on
average, consistent with ROC-AUC saturating at the $+$T level rather
than with community features being harmful.

\textbf{Logistic Regression shows no community gain on any metric.}
This is consistent with the community signal being non-linear and
inaccessible to a linear model on these features.

At the descriptive level, the data are not consistent with the
naive ``distrust at community boundaries'' reading. Aggregating
across all rolling-fold training blocks, same-community pairs have
an empirical positive rate of $0.0140\%$ ($215 / 1{,}540{,}969$)
while cross-community pairs are at $0.0095\%$ ($330 / 3{,}472{,}377$);
same-community pairs are slightly \emph{more} likely to receive a
new negative rating in the next window (Fig.~\ref{fig:samecross}).

\begin{figure}[t]
\centering
\includegraphics[width=0.85\columnwidth]{results/figures/same_vs_cross_positive_rate.png}
\caption{Aggregate training positive rate by community membership.
Same-community pairs are $\sim 47\%$ more likely to receive a new
negative rating than cross-community pairs. The naive
``cross-community distrust'' hypothesis is not supported.}
\label{fig:samecross}
\end{figure}

In the fitted Logistic Regression on S$+$T$+$C,
\code{neighborhood\_overlap} is the strongest single community
feature (standardized coefficient $-1.36$, in the direction that
lower shared-neighborhood predicts higher new-negative probability);
it is also among the top-ranked community features in Random Forest
(Tables~\ref{tab:imp_lr_topfam} and \ref{tab:imp_rf_topfam} in
Section~\ref{sec:results:locked}). The coarse same/cross indicator
carries little weight once \code{neighborhood\_overlap}, boundary
fractions, and community sizes are present.

\textbf{Precision@$k$ is unstable under rolling.} Per-fold median
P@50 is zero for every (model, feature set), and median P@100 is at
most $0.01$ for the best (model, feature set) combination. Most
folds return no true positives in the top 50. We therefore demote
Precision@$k$ to a supplementary metric.

\subsection{Comparison with stronger baselines}
\label{sec:results:baselines}

To test whether the temporal lift is genuinely broader than a
trivial recency effect, and to bound the headroom for a stronger
tabular learner, we add (i) four parameter-free heuristic baselines
and (ii) LightGBM with a single fixed configuration, evaluated under
the same rolling-origin protocol of Section~\ref{sec:setup:rolling}.

\textbf{Recency-only ranker.} The recency-only heuristic is a
within-fold rank-sum
$\sum_j \mathrm{sign}_j \cdot \mathrm{rank}_j(x_j)$ over six recency
features (days-since-last-activity for $u$ and $v$; recent and
exponentially-decayed negative counts on both sides). Signs are
hand-chosen for intuitive directionality; no parameters are learned.
Across the 25 folds it reaches ROC-AUC
$0.886 \pm 0.062$, statistically indistinguishable from
structural-only Logistic Regression ($0.889 \pm 0.077$) and Random
Forest ($0.881 \pm 0.073$). The learned $\mathrm{S}{+}\mathrm{T}$
models still beat this heuristic clearly (RF $0.924$, LR $0.930$),
but the headroom over a parameter-free recency baseline is
$\sim 0.04$ ROC-AUC rather than the $\sim 0.07$ visible above the
structural-only baseline. \emph{Much of what we labeled ``temporal
lift'' in Section~\ref{sec:results:temporal} is recency lift}; the
engineered temporal feature family adds a real but modest increment
above pure recency.

\textbf{Victim-, issuer-, and structural-only rank-sums.}
Single-sided risk summaries do not suffice for the pair-level task.
Victim-risk-only reaches ROC-AUC $0.698 \pm 0.109$, issuer-risk-only
$0.755 \pm 0.127$, and a structural-only rank-sum (CN, Jaccard, AA,
PA, reciprocity, prior interaction) $0.772 \pm 0.073$, all
substantially below the learned models. Pair-level prediction
requires combining both sides; one-sided summaries alone are not
informative enough to compete.

\textbf{LightGBM at a fixed untuned configuration does not
outperform Random Forest on this task.} We use
\code{n\_estimators=200}, \code{learning\_rate=0.05},
\code{num\_leaves=31}, \code{min\_child\_samples=20}, and
\code{scale\_pos\_weight}$ = n_-/n_+$; no hyperparameter search.
LightGBM ROC-AUC is significantly below Random Forest on every
feature set (Wilcoxon $p \leq 0.002$ under BH-FDR $q = 0.05$). On
PR-AUC the comparison is mixed: LightGBM is slightly better on S
and S$+$T but slightly worse on S$+$T$+$C. None of the four
pre-declared LightGBM-centric primary tests
(LGBM (S+T) - LGBM S on ROC-AUC and PR-AUC, LGBM (S+T+C) - LGBM
(S+T) on PR-AUC, LGBM (S+T) - RF (S+T) on PR-AUC) survive
Holm-Bonferroni; the temporal ROC-AUC lift under LightGBM is
directionally consistent (17 of 25 folds positive) but does not
reach the corrected threshold ($p = 0.017$, Holm threshold
$0.0125$). We do not claim that LightGBM is fundamentally worse than
Random Forest on this task; a more aggressive tuning could change the
picture. We report the comparison to document that the interpretable
Random Forest result is not obviously leaving accuracy on the table.

\textbf{The community effect is not robust across model families.}
Under Random Forest the (S$+$T$+$C) - (S$+$T) PR-AUC delta is small,
positive, and marginally significant (Section~\ref{sec:results:community});
under LightGBM the same delta has a slightly negative point estimate
(mean $-0.007$, 14/25 folds positive, Wilcoxon $p = 0.51$). The
(S$+$T$+$C) - (S$+$T) ROC-AUC delta is directionally positive under
LightGBM ($+0.009$, 19/25 folds positive, $p = 0.026$) but null
under Random Forest. The marginal Phase-D significance of the RF
PR-AUC community effect therefore should not be read as a universal
claim about community-aware features; the direction is
model-dependent.

\subsection{Held-out single-fold analysis (secondary)}
\label{sec:results:locked}

For continuity with prior pipeline stages,
Table~\ref{tab:test_metrics} reports the locked single held-out fold
metrics. The locked fold's RF (S$+$T$+$C) PR-AUC ($0.0703$) is well
above the rolling-fold mean ($0.017$); the locked fold thus
over-represents the magnitude of the community PR-AUC effect that is
visible across all 25 folds. We retain the table, the curve
overlays, the importance rankings, and the bootstrap analysis here
because they remain informative about model behavior on a single
representative fold, but the headline numerical claims in this paper
come from the rolling-origin protocol of
Section~\ref{sec:results:rolling}.

\begin{table}[t]
\caption{Held-out single-fold metrics ($t \in \{48, 49\}$;
34 positives, 148{,}639 candidate pairs). Secondary to
Table~\ref{tab:rolling_summary}.}
\label{tab:test_metrics}
\centering
\begin{tabular}{llcccc}
\toprule
F.\,set & Model & ROC-AUC & PR-AUC & P@50 & P@100 \\
\midrule
S      & logreg & 0.909 & 0.0034 & 0.00 & 0.01 \\
S$+$T   & logreg & 0.949 & 0.0059 & 0.00 & 0.01 \\
S$+$T$+$C & logreg & 0.945 & 0.0058 & 0.00 & 0.01 \\
\addlinespace
S      & rf     & 0.880 & 0.0223 & 0.04 & 0.03 \\
S$+$T   & rf     & 0.950 & 0.0400 & 0.08 & 0.07 \\
S$+$T$+$C & rf     & 0.940 & 0.0703 & 0.10 & 0.06 \\
\bottomrule
\end{tabular}
\end{table}

\subsubsection{ROC and PR curves on the locked fold}
Fig.~\ref{fig:rf_roc} and Fig.~\ref{fig:rf_pr} overlay Random Forest
ROC and PR curves for the three feature sets on the locked fold.
The S$+$T (blue) and S$+$T$+$C (red) ROC curves nearly coincide in
the operationally relevant low-FPR region; both are clearly above
the structural-only S (grey) curve, visualizing the ROC saturation
that the rolling-origin tests also report. The S$+$T$+$C PR curve
sits visibly above S$+$T between recall $\sim 0.05$ and $\sim 0.30$
on this fold. As noted in Table~\ref{tab:rolling_summary}, the
average across-fold separation is much smaller; this locked fold is
on the favorable end of the per-fold PR-AUC distribution.

\begin{figure}[t]
\centering
\includegraphics[width=\columnwidth]{results/figures/test_rf_roc.png}
\caption{Test ROC curves for Random Forest on the locked single fold.
S$+$T and S$+$T$+$C are visually indistinguishable in the
operationally relevant low-FPR region; both clearly dominate the
structural-only S curve. Curve shapes generalize qualitatively;
absolute AUC values on this single fold are above the rolling-origin
average reported in Table~\ref{tab:rolling_summary}.}
\label{fig:rf_roc}
\end{figure}

\begin{figure}[t]
\centering
\includegraphics[width=\columnwidth]{results/figures/test_rf_pr.png}
\caption{Test precision-recall curves for Random Forest on the
locked single fold, log precision axis. The community-augmented
S$+$T$+$C curve sits visibly above S$+$T between recall $\sim 0.05$
and $\sim 0.30$ on this fold; the rolling-origin average separation
is substantially smaller (Table~\ref{tab:rolling_summary}). Dashed
line marks the base rate of $2.29 \times 10^{-4}$.}
\label{fig:rf_pr}
\end{figure}

Logistic Regression overlays (not shown; available on disk as
\code{test\_logreg\_\{roc,pr\}.png}) tell a different story: the
S$+$T and S$+$T$+$C curves are essentially indistinguishable, while
S is below. This visually confirms the absence of any LR community
gain.

\subsubsection{Feature importance (qualitative)}
Tables~\ref{tab:imp_lr_topfam} and~\ref{tab:imp_rf_topfam} list the
top-5 features per family for each model on the S$+$T$+$C locked-fold
fit. Coefficient magnitudes (LR) and Gini importances (RF) are
qualitative diagnostics; the inferential conclusions are carried by
the rolling-origin tests.

\begin{table}[t]
\caption{Top-5 features per family, Logistic Regression on S$+$T$+$C
(standardized coefficients; $+$ drives the positive class). Locked-fold
fit; qualitative diagnostic.}
\label{tab:imp_lr_topfam}
\centering
\small
\begin{tabular}{llr}
\toprule
Family & Feature & Coef. \\
\midrule
struct.\ & cn                          & $+1.28$ \\
struct.\ & prior\_interaction          & $-1.17$ \\
struct.\ & pos\_cn                     & $-1.09$ \\
struct.\ & reciprocity\_flag           & $+0.88$ \\
struct.\ & jaccard                     & $+0.88$ \\
\addlinespace
temp.\   & days\_since\_last\_activity\_u  & $-1.29$ \\
temp.\   & days\_since\_last\_uv\_interaction & $+1.19$ \\
temp.\   & uv\_interaction\_count\_past & $-1.17$ \\
temp.\   & decayed\_total\_activity\_v & $+0.76$ \\
temp.\   & decayed\_total\_activity\_u & $+0.71$ \\
\addlinespace
comm.\   & neighborhood\_overlap       & $-1.36$ \\
comm.\   & boundary\_frac\_u           & $+0.17$ \\
comm.\   & community\_size\_diff\_log  & $-0.10$ \\
comm.\   & community\_size\_v          & $-0.09$ \\
comm.\   & same\_community             & $+0.08$ \\
\bottomrule
\end{tabular}
\end{table}

\begin{table}[t]
\caption{Top-5 features per family, Random Forest on S$+$T$+$C
(Gini importance). Locked-fold fit; qualitative diagnostic.}
\label{tab:imp_rf_topfam}
\centering
\small
\begin{tabular}{llr}
\toprule
Family & Feature & Importance \\
\midrule
struct.\ & aa                          & $0.050$ \\
struct.\ & pa                          & $0.037$ \\
struct.\ & neg\_in\_v                  & $0.034$ \\
struct.\ & reciprocity\_flag           & $0.032$ \\
struct.\ & neg\_received\_ratio\_v     & $0.030$ \\
\addlinespace
temp.\   & decayed\_total\_activity\_u & $0.109$ \\
temp.\   & days\_since\_last\_activity\_u & $0.095$ \\
temp.\   & recent\_total\_given\_u     & $0.053$ \\
temp.\   & decayed\_neg\_received\_v   & $0.051$ \\
temp.\   & decayed\_neg\_given\_u      & $0.051$ \\
\addlinespace
comm.\   & boundary\_frac\_v           & $0.018$ \\
comm.\   & neighborhood\_overlap       & $0.018$ \\
comm.\   & boundary\_frac\_u           & $0.017$ \\
comm.\   & community\_size\_v          & $0.017$ \\
comm.\   & community\_size\_u          & $0.016$ \\
\bottomrule
\end{tabular}
\end{table}

\subsubsection{Bootstrap uncertainty on the locked fold}
\label{sec:results:bootstrap}

Tables~\ref{tab:boot_abs} and~\ref{tab:boot_delta} summarize the
paired stratified bootstrap ($B = 1000$) on the locked single fold.
These CIs quantify sampling uncertainty conditional on this fold;
they are complementary to the rolling-origin Wilcoxon tests rather
than primary. Note that the locked-fold community PR-AUC effect
($+0.026$, 94\% of bootstrap resamples positive, CI
$[-0.003, +0.068]$) is consistent in direction with the rolling
result but several times larger than the across-fold mean of $+0.003$.

\begin{table}[t]
\caption{Bootstrap 95\% percentile CIs on absolute test metrics
(Random Forest, locked single fold). Secondary to
Table~\ref{tab:rolling_summary}.}
\label{tab:boot_abs}
\centering
\small
\begin{tabular}{lcc}
\toprule
F.\,set & ROC-AUC & PR-AUC \\
\midrule
S        & $0.88$ $[0.80, 0.94]$ & $0.030$ $[0.003, 0.097]$ \\
S$+$T    & $0.95$ $[0.92, 0.97]$ & $0.052$ $[0.011, 0.131]$ \\
S$+$T$+$C & $0.94$ $[0.90, 0.97]$ & $0.079$ $[0.012, 0.180]$ \\
\midrule
F.\,set & P@50 & P@100 \\
\midrule
S        & $0.05$ $[0.00, 0.12]$ & $0.03$ $[0.00, 0.07]$ \\
S$+$T    & $0.08$ $[0.02, 0.16]$ & $0.07$ $[0.02, 0.11]$ \\
S$+$T$+$C & $0.09$ $[0.02, 0.18]$ & $0.06$ $[0.03, 0.11]$ \\
\bottomrule
\end{tabular}
\end{table}

\begin{table}[t]
\caption{Bootstrap 95\% percentile CIs on paired metric deltas
(Random Forest, locked single fold). \emph{frac\_boot} is the share
of resamples with $\Delta > 0$; it is not a p-value.}
\label{tab:boot_delta}
\centering
\small
\begin{tabular}{lllc}
\toprule
$\Delta$ & Metric & 95\% CI & frac\_boot \\
\midrule
$(\mathrm{S}{+}\mathrm{T}) - \mathrm{S}$
  & ROC-AUC & $[+0.028, +0.123]$ & $1.000$ \\
  & PR-AUC  & $[-0.021, +0.082]$ & $0.849$ \\
  & P@50    & $[-0.040, +0.100]$ & $0.755$ \\
  & P@100   & $[-0.000, +0.080]$ & $0.915$ \\
\addlinespace
$(\mathrm{S}{+}\mathrm{T}{+}\mathrm{C}) - (\mathrm{S}{+}\mathrm{T})$
  & ROC-AUC & $[-0.027, +0.002]$ & $0.073$ \\
  & PR-AUC  & $[-0.003, +0.068]$ & $0.935$ \\
  & P@50    & $[+0.000, +0.060]$ & $0.465$ \\
  & P@100   & $[-0.030, +0.020]$ & $0.157$ \\
\addlinespace
$(\mathrm{S}{+}\mathrm{T}{+}\mathrm{C}) - \mathrm{S}$
  & ROC-AUC & $[+0.015, +0.113]$ & $0.998$ \\
  & PR-AUC  & $[+0.006, +0.120]$ & $0.999$ \\
  & P@50    & $[-0.020, +0.120]$ & $0.850$ \\
  & P@100   & $[-0.010, +0.080]$ & $0.907$ \\
\bottomrule
\end{tabular}
\end{table}

% =========================================================
\section{Error Analysis and Interpretation}
% =========================================================

We separate two kinds of evidence about community-level signal. The
first is \emph{global}, drawn from feature distributions and fitted
feature importances over aggregate training data. The second is
\emph{local}, drawn from inspection of the top-ranked test pairs from
the best locked-fold model. The local analysis is descriptive and
based on a very small number of true positives; we read it as a
model-behavior diagnostic, not as confirmation or refutation of any
sociological hypothesis.

\subsection{Global, feature-level evidence}

Two descriptive facts at the population level are worth stating
separately. They are aggregated across all 25 rolling-fold training
blocks; the same direction holds for any individual fold's training
block.

\begin{enumerate}
\item Same-community pairs have an empirical positive rate of
$0.0140\%$ ($215 / 1{,}540{,}969$), whereas cross-community pairs are
at $0.0095\%$ ($330 / 3{,}472{,}377$). The naive ``cross-community
distrust'' reading is not supported on this dataset
(Fig.~\ref{fig:samecross}); same-community pairs are slightly
\emph{more} likely to receive a new negative rating in the next month.
\item In the fitted Logistic Regression on S$+$T$+$C,
\code{neighborhood\_overlap} is the strongest pair-level community
feature (standardized coefficient $-1.36$), in the direction that
lower shared-neighborhood predicts higher new-negative probability.
The same feature is among the top community-family signals in Random
Forest, where the dominant community-family signals are the boundary
fractions.
\end{enumerate}

Taken together, these two facts are consistent with -- but do not by
themselves prove -- a refined ``intra-community weak tie'' reading:
distrust does not predominantly appear at coarse community boundaries,
yet pair-level neighborhood overlap still matters and points in the
Granovetter direction. Establishing this reading more firmly would
require conditional analyses (e.g., overlap-conditioned positive rates
inside and across communities), which we leave to future work.

\subsection{Local diagnostic from the top-50 scored pairs (locked
fold, RF on S$+$T$+$C)}

At the top-of-ranked-list operating point ($k = 50$), the best model
returns 5 true positives and 45 false positives. With $n_{TP} = 5$ the
sample is too small for statistical comparison; what follows is a
qualitative look at the kinds of pairs the model surfaces.

\begin{table}[t]
\caption{Mean feature values among the top-50 RF S$+$T$+$C scored
test pairs on the locked fold (illustrative; $n_{TP} = 5$, $n_{FP} = 45$).}
\label{tab:top50}
\centering
\small
\begin{tabular}{lrr}
\toprule
Feature & TPs & FPs \\
\midrule
bootstrap proba              & $0.941$ & $0.914$ \\
neg\_out\_u                  & $17.80$ & $31.53$ \\
neg\_in\_v                   & $25.60$ & $14.56$ \\
decayed\_neg\_received\_v    & $1.34$  & $2.04$  \\
days\_since\_last\_activity\_u & $3.59$ & $2.53$ \\
same\_community              & $0.00$  & $0.16$  \\
neighborhood\_overlap        & $0.080$ & $0.055$ \\
boundary\_frac\_u            & $0.295$ & $0.394$ \\
uv\_interaction\_count\_past & $0.00$  & $0.00$  \\
\bottomrule
\end{tabular}
\end{table}

Three qualitative patterns are visible
(Table~\ref{tab:top50}). First, the model concentrates risk on a
small set of recurring ``issuer'' nodes (high \code{neg\_out\_u}) and
recurring ``victim'' nodes (high \code{neg\_in\_v}); this matches the
features. Second, what distinguishes the small true-positive subset
from the false positives at the very top is mostly on the
\emph{victim} side (mean \code{neg\_in\_v} 25.6 vs 14.6 for FPs), not
the issuer side. False positives at the top tend to be highly active
negative-givers whose specific $(u, v)$ target did not receive a new
negative in the test window. Third, on the community-aware features
the local direction in this slice does \emph{not} mirror the global
TRAIN pattern: in the top-50, all 5 TPs are cross-community, and the
mean \code{neighborhood\_overlap} is higher (not lower) for TPs than
for FPs.

We deliberately do not interpret the top-50 slice as confirming or
refuting the global pattern. With 5 true positives, the local
direction is sensitive to which specific candidate pairs the model
happened to rank highly, and the local statistic does not have the
sample support to override the population-level evidence from the
previous subsection. Supplementary top-200 and top-500 slices
(available on disk as \code{results/stage7\_error\_top200.csv} and
\code{results/stage7\_error\_top500.csv}) further illustrate that
the model is highly \emph{front-loaded}: only 7 of the top 200 and
8 of the top 500 scored pairs are true positives. Most of the model's
correct mass is concentrated in a small number of high-confidence
recurring patterns.

\subsection{Summary}

The descriptive evidence on training data is consistent with a
refined weak-tie reading in which lower pair-level neighborhood
overlap is associated with higher distrust risk, while cross-community
status by itself is not. The local top-50 diagnostic on the locked
fold shows that the model's highest-confidence predictions are
concentrated on recurring victims and that its top false positives
are dominated by highly active negative-givers; the community-feature
direction in this small slice does not align with the global
feature-importance direction, which may reflect a conditional-versus-marginal
divergence, although the sample is too small for a strong conclusion.

% =========================================================
\section{Limitations}\label{sec:lim}
% =========================================================

We list the limitations explicitly so that follow-up work can
prioritize.

\begin{itemize}
\item \textbf{Per-fold positive counts are still small.} Although
the rolling-origin protocol accumulates 425 positive test events
across 25 folds, individual folds have only $\sim 17$ positives on
average, and per-fold PR-AUC has standard deviation roughly twice its
mean for Random Forest. This is why the community PR-AUC effect sits
at the margin of statistical significance even with $n = 25$.
\item \textbf{Single network, no external replication.} We do not
report results on Bitcoin-Alpha or any other signed network.
Bitcoin-Alpha is a natural transfer target and the most important
next step for reviewer credibility.
\item \textbf{Much of the temporal lift is recency.} A parameter-free
recency-only ranker (Section~\ref{sec:results:baselines}) already
reaches ROC-AUC $0.886$, comparable to structural-only learned models.
The engineered temporal feature family adds a real but modest
increment above raw recency ($\sim 0.04$ ROC-AUC); the qualitative
story should be read as ``recency information dominates'' rather than
``our specific decay / recent-window feature engineering does.''
\item \textbf{Limited baseline scope.} Beyond Logistic Regression,
Random Forest, and LightGBM under a fixed untuned configuration, we
do not benchmark more aggressive tabular models or graph embeddings.
LightGBM did not outperform Random Forest at our configuration; a
tuned variant could.
\item \textbf{Community effect is not robust across model families.}
The marginal RF PR-AUC community gain does not replicate under
LightGBM (mean $-0.007$, $p = 0.51$ on that comparison). Conclusions
about community-aware features are RF-specific in our experiments.
\item \textbf{Marginal nature of the Holm-surviving RF community
PR-AUC test.} The primary-family Wilcoxon $p = 0.048$ sits at the
loosest Holm threshold. A different primary family, a stricter
correction (e.g.\ Bonferroni on all 24 tests at
$\alpha / 24 = 0.002$), or a small data perturbation could flip this
decision. We report it honestly rather than soft-pedalling it.
\item \textbf{Precision@$k$ is unstable under the rolling protocol.}
Per-fold median P@50 is zero for every (model, feature set). We
report P@$k$ only as supplementary evidence and do not use it to
support community claims.
\item \textbf{Last-rating-wins snapshot collapse.} The snapshot
graph uses the latest rating per pair; this represents current trust
judgment but discards multi-event nuance. We restore the full event
history only for temporal and trust-summary features.
\item \textbf{Bootstrap CIs apply to the locked single fold.}
Section~\ref{sec:setup:bootstrap}'s percentile CIs quantify sampling
uncertainty conditional on the locked fold and can under-cover when
the bootstrap distribution is skewed (e.g.\ for Precision@$k$ near
its hard ceiling). The rolling-origin Wilcoxon results are our
primary uncertainty story.
\item \textbf{\code{prior\_interaction} is partly degenerate.} By
construction of the candidate set, this flag is $1$ only when $u$
previously rated $v$ positively. We retain it for transparency but
note that its fitted coefficient should be read with care.
\item \textbf{\emph{frac\_boot}$(>0)$ is not a p-value.} We use it
only to summarize the direction of the bootstrap distribution.
\end{itemize}

% =========================================================
\section{Conclusion}
% =========================================================

We presented an interpretable, leakage-controlled pipeline for early
warning of new negative trust-edge emergence in the Bitcoin-OTC
temporal signed network, together with a three-way feature ablation
under Logistic Regression and Random Forest. We evaluated under a
rolling-origin protocol over 25 disjoint single-month test folds
(425 total positive events), with paired Wilcoxon and sign tests on
per-fold metric deltas and a pre-declared primary family controlled
by Holm-Bonferroni, and compared the pipeline against parameter-free
heuristic baselines and LightGBM.

\textbf{Temporal information robustly improves early warning.} The
$(\mathrm{S}{+}\mathrm{T}) - \mathrm{S}$ ROC-AUC delta is positive
in 22 of 25 RF folds and 21 of 25 LR folds, with Wilcoxon p-values
$\leq 1.4\!\times\!10^{-4}$ under both models after Holm correction
and matched-pairs rank-biserial $r_{\mathrm{rb}} \approx +0.8$. The
Random Forest temporal PR-AUC lift is also significant after Holm
correction ($p = 6.3\!\times\!10^{-4}$,
$r_{\mathrm{rb}} = +0.74$).

\textbf{A large fraction of the temporal lift is recency.} A
parameter-free recency-only rank-sum reaches ROC-AUC $0.886$ across
the same folds, indistinguishable from structural-only learned
models. The engineered temporal feature family still beats it
clearly under $\mathrm{S}{+}\mathrm{T}$ (RF $0.924$, LR $0.930$), but
much of the ``temporal'' improvement we report should be understood
as recency information rather than as specific evidence for any
particular temporal feature design.

\textbf{Community-aware features provide a small, PR-oriented gain
under Random Forest, but the effect is model-dependent.} The
$(\mathrm{S}{+}\mathrm{T}{+}\mathrm{C}) - (\mathrm{S}{+}\mathrm{T})$
PR-AUC delta is positive in 19 of 25 RF folds (sign-test
$p = 0.015$; Wilcoxon $p = 0.048$, marginal under the pre-declared
Holm threshold; $r_{\mathrm{rb}} = +0.45$). Under LightGBM the same
delta is slightly negative on average. We frame the community PR-AUC
effect as small, RF-specific, and marginally significant rather than
as a general claim.

\textbf{Community features do not improve ROC-AUC over the temporal
baseline.} Neither RF nor LR shows a significant community ROC-AUC
gain ($p = 0.096$ and $p = 0.47$ respectively), and the RF
direction is mildly negative on average -- consistent with saturation.

\textbf{Logistic Regression shows no community gain on any metric},
consistent with the community signal being non-linear.

\textbf{LightGBM at a fixed untuned configuration does not outperform
Random Forest} on this task. We do not claim it cannot; we report
what a defensible default configuration produces.

\textbf{Precision@$k$ is unstable under the rolling protocol} --
most folds return no true positives in the top 50 -- and we
therefore demote it from headline status to a supplementary metric.

At the descriptive level, the data are not consistent with a coarse
``distrust between communities'' reading: same-community pairs in
aggregate training data have a higher empirical positive rate than
cross-community pairs. They are consistent with a finer pair-level
pattern in which lower neighborhood overlap is associated with
higher new-negative probability. We frame this as a refined,
conservative reading rather than as a confirmed weak-tie result.

The largest natural extensions are: (i) external replication on
Bitcoin-Alpha, the single highest-value next step; (ii) more
aggressive baselines, including tuned gradient-boosted trees and
shallow graph embeddings; (iii) a richer conditional analysis of
community signal (e.g.\ overlap-conditioned positive rates inside
and across communities); and (iv) BCa or studentized bootstrap CIs
for the locked-fold analysis to address skew at the Precision@$k$
ceiling.

% =========================================================
% REFERENCES
% =========================================================
\begin{thebibliography}{99}

\bibitem{kumar2016}
S.~Kumar, F.~Spezzano, V.~S. Subrahmanian, and C.~Faloutsos,
``Edge weight prediction in weighted signed networks,''
in \emph{Proc.\ IEEE 16th Int.\ Conf.\ Data Mining (ICDM)},
Barcelona, Spain, Dec.\ 2016, pp.\ 221--230.
DOI: 10.1109/ICDM.2016.0033.

\bibitem{bertazzi2018}
I.~Bertazzi, S.~Huet, G.~Deffuant, and F.~Gargiulo,
``The anatomy of a web of trust: The Bitcoin-OTC market,''
in \emph{Social Informatics (SocInfo 2018)},
Lecture Notes in Computer Science, vol.\ 11185, Springer, 2018,
pp.\ 228--241.
DOI: 10.1007/978-3-030-01129-1\_14.

\bibitem{choudhury2024}
N.~Choudhury,
``Community-aware evolution similarity for link prediction in dynamic
social networks,''
\emph{Mathematics (MDPI)}, vol.\ 12, no.\ 2, art.\ 285, 2024.
DOI: 10.3390/math12020285.

\bibitem{zhang2024}
Z.~Zhang, P.~Zhao, X.~Li, J.~Liu, X.~Zhang, J.~Huang, and X.~Zhu,
``Signed graph representation learning: A survey,''
\emph{arXiv preprint arXiv:2402.15980}, Feb.\ 2024.

\bibitem{granovetter1973}
M.~S. Granovetter,
``The strength of weak ties,''
\emph{Amer.\ J.\ Sociology}, vol.\ 78, no.\ 6, pp.\ 1360--1380, 1973.
DOI: 10.1086/225469.

\bibitem{blondel2008}
V.~D. Blondel, J.-L. Guillaume, R.~Lambiotte, and E.~Lefebvre,
``Fast unfolding of communities in large networks,''
\emph{J.\ Stat.\ Mech.: Theory Exp.}, vol.\ 2008, no.\ 10, art.\ P10008,
Oct.\ 2008. DOI: 10.1088/1742-5468/2008/10/P10008.

\bibitem{adamic2003}
L.~A. Adamic and E.~Adar,
``Friends and neighbors on the web,''
\emph{Social Networks}, vol.\ 25, no.\ 3, pp.\ 211--230, 2003.
DOI: 10.1016/S0378-8733(03)00009-1.


\bibitem{snap}
J.~Leskovec and A.~Krevl,
``SNAP datasets: Stanford large network dataset collection,''
\url{http://snap.stanford.edu/data}, Jun.\ 2014.
Bitcoin-OTC subset:
\url{https://snap.stanford.edu/data/soc-sign-bitcoinotc.html}.

\end{thebibliography}

\end{document}
